\chapter{Kết Luận}

\newpage

\section{Tóm tắt kết quả đạt được}

Phần 2 của đồ án đã triển khai thành công một pipeline phân tích dữ liệu hoàn chỉnh—từ
khảo sát dữ liệu (EDA), tiền xử lý, xây dựng mô hình đến đánh giá kết quả—trên bộ
dữ liệu ảnh cộng hưởng từ não bộ theo dõi dọc \textbf{OASIS-2} (373 lượt khám của
150 bệnh nhân, nhóm tuổi 60 - 96).
 
\subsection{ Pipeline Tiền Xử Lý}
 
\begin{itemize}
  \item \textbf{Xử lý missing values có căn cứ thống kê:}
  
   Biến \texttt{SES}
        (tỉ lệ khuyết \(\approx 5{,}09\%\)) được impute bằng
        \emph{Group-by-Median} theo nhóm học vấn \texttt{EDUC}, tận dụng tương
        quan thứ bậc giữa hai biến thay vì dùng median toàn cục. Hai quan sát
        khuyết ở biến mục tiêu \texttt{MMSE} bị loại bỏ (listwise deletion) để
        tránh sai lệch ước lượng.
 
  \item \textbf{Phát hiện và xử lý đa cộng tuyến:} 
  
  Phân tích Centered VIF xác
        nhận bốn đặc trưng y sinh (\texttt{nWBV}, \texttt{Age}, \texttt{EDUC},
        \texttt{SES}) hoàn toàn độc lập tuyến tính (\(\text{VIF} < 3\)) sau khi
        tích hợp hệ số chặn (intercept). Biến \texttt{EDUC} bị loại bỏ ở mô
        hình OLS chọn biến do mất ý nghĩa thống kê sau hiệu chỉnh sai số chuẩn
        mạnh HC3 (\(p > 0{,}05\)).
 
  \item \textbf{Nguyên tắc fit-on-train, transform-both:}
  
   Toàn bộ tham số tiền
        xử lý (giá trị median theo nhóm, tham số chuẩn hóa Z-score) được ước
        lượng \emph{hoàn toàn} từ tập huấn luyện và áp dụng thụ động lên tập
        kiểm thử, đảm bảo không rò rỉ thông tin (data leakage).
\end{itemize}
 
\subsection{So Sánh Bốn Mô Hình trên Tập Kiểm Thử}
 
Bốn kiến trúc mô hình được xây dựng, huấn luyện và đánh giá độc lập trên tập
kiểm thử (\(n_{\text{test}} = 75\) bệnh nhân chưa từng thấy):
 
\begin{table}[h]
  \centering
  \caption{Kết quả đánh giá trên tập kiểm thử (Test Set)}
  \label{tab:model_comparison}
  \begin{tabular}{lrrr}
    \toprule
    \textbf{Mô hình} & \textbf{MAE} & \textbf{RMSE} & \boldmath$R^2_{\text{Test}}$ \\
    \midrule
    OLS Baseline                                     & 2{,}4563 & 2{,}9701 & $-0{,}1338$ \\
    OLS Chọn biến (Centered VIF)                     & 2{,}4590 & 2{,}9698 & $-0{,}1337$ \\
    \textbf{Ridge} $(\lambda^* = 6{,}5793)$          & \textbf{2{,}2899} & \textbf{2{,}9432} & $\mathbf{-0{,}1106}$ \\
    Kernel Ridge $(\alpha^*,\,\gamma^*$ via CV$)$    & 2{,}19-2{,}69 & 2{,}88-3{,}51 & $-0{,}10$ đến $-0{,}42$ \\
    \bottomrule
  \end{tabular}
\end{table}
 
\textbf{Ridge Regression} $(\lambda^* = 6{,}5793$, chọn qua K-Fold CV với
\(k=5\)) là mô hình có hiệu năng tổng quát hóa tốt nhất trên cả ba chỉ số.
Cơ chế phạt $L_2$ phát huy hiệu quả kép: thu nhỏ đồng đều hệ số của các biến
có tương quan nội tại (\texttt{Age}--\texttt{nWBV}, \texttt{EDUC}--\texttt{SES})
đồng thời giữ nguyên toàn bộ không gian đặc trưng—khác biệt về triết lý so với
\emph{feature selection} của OLS chọn biến.
 
\textbf{Kernel Ridge Regression} vượt trội về diagnostic plots trên tập huấn
luyện (phần dư phân tán đều, Scale-Location phẳng hơn), song thua Ridge trên
tập kiểm thử kể cả sau khi $(\alpha^*, \gamma^*)$ được tối ưu qua CV Grid
Search $(6 \times 5 = 30$ tổ hợp$)$. Kết quả này xác nhận \emph{nghịch lý kích
thước mẫu}: với $n = 371$ và $p = 4$, ma trận Gram
$K \in \mathbb{R}^{296 \times 296}$ có bậc tự do vượt xa Ridge (chỉ học
$4$ hệ số $\hat{\beta}$), dẫn đến overfit cấu trúc Train mà không tổng quát
hóa được.
 
\subsection{Ý Nghĩa Lâm Sàng của Các Đặc Trưng}
 
\begin{itemize}
  \item \textbf{Thể tích não chuẩn hóa (\texttt{nWBV}):} 
  
  Đặc trưng mang tính
        sinh học trực tiếp nhất. Sự sụt giảm \texttt{nWBV} phản ánh quá trình
        teo tế bào thần kinh do tích tụ mảng bám amyloid và nút thắt sợi
        tau—dấu ấn bệnh lý đặc trưng của Alzheimer. Mô hình ghi nhận tương quan
        thuận mạnh giữa tỷ lệ thể tích chất xám/chất trắng còn lại và khả năng
        duy trì điểm số nhận thức MMSE.
 
  \item \textbf{Tuổi tác (\texttt{Age}):} 
  
    Yếu tố nguy cơ không thể can thiệp
        hàng đầu. Lão hóa thần kinh tích lũy theo thời gian tạo ra mối quan hệ
        phi tuyến với MMSE—đặc biệt ở giai đoạn bùng phát triệu chứng lâm sàng
        muộn.
 
  \item \textbf{Học vấn (\texttt{EDUC}) và Kinh tế--Xã hội (\texttt{SES}):}

        Hai yếu tố cấu thành \emph{Dự trữ nhận thức} (Cognitive Reserve). Bệnh
        nhân có nền tảng học vấn cao duy trì mạng lưới kết nối thần kinh dự
        phòng tốt hơn, giúp trì hoãn biểu hiện lâm sàng dù cấu trúc não đã bắt
        đầu tổn thương.
\end{itemize}

\newpage

\section{Khó khăn gặp phải}

\begin{itemize}
  \item \textbf{Kích thước mẫu hạn chế và cấu trúc dọc (Longitudinal):}

        Bộ dữ liệu OASIS-2 chỉ có 150 bệnh nhân độc lập với 373 lượt khám.
        Việc phân chia ngẫu nhiên các lượt khám có nguy cơ để cùng một bệnh
        nhân xuất hiện ở cả tập huấn luyện lẫn tập kiểm thử (data leakage
        theo chiều thời gian). Ngoài ra, kích thước mẫu nhỏ khiến Kernel Ridge
        không thể tổng quát hóa: ma trận Gram $K \in \mathbb{R}^{296 \times
        296}$ gần như ghi nhớ toàn bộ tập Train mà không học được quy luật
        tổng quát.
 
  \item \textbf{Vi phạm giả định tuyến tính và phương sai thay đổi:}

        Tất cả bốn mô hình đều ghi nhận $R^2 < 0$ và diagnostic plots cho thấy
        rõ ràng hai vi phạm giả định Gauss--Markov: (i)~phần dư có cấu trúc
        hình phễu/đường cong (tính phi tuyến), và (ii)~đường Scale-Location có
        độ dốc (phương sai thay đổi). Các vi phạm này phản ánh bản chất phi
        tuyến của quá trình suy giảm nhận thức sinh học, không phải lỗi kỹ
        thuật của pipeline.
 
  \item \textbf{Đa cộng tuyến ảo do thiếu intercept:}

        Phân tích VIF ban đầu (Uncentered) báo cáo giá trị lạm phát cao giả tạo
        do hệ số chặn bị tính vào phương trình. Việc chuyển sang Centered VIF
        (tích hợp intercept trước khi tính) giải quyết vấn đề và xác nhận bốn
        đặc trưng thực sự độc lập tuyến tính (\(\text{VIF} < 3\)).
 
  \item \textbf{Tối ưu hóa siêu tham số không phản ánh đúng khả năng tổng
        quát hóa:}

        K-Fold CV thông thường với \(k = 5\) có thể chia cùng một bệnh nhân vào
        cả fold Train lẫn fold Validation, làm CV MSE lạc quan hơn thực tế.
        Ngoài ra, Kernel Ridge đạt CV MSE tốt trên tập Train nhưng vẫn thua
        Ridge trên tập Test—bộc lộ khoảng cách giữa \emph{validation
        performance} và \emph{generalization performance} khi kích thước mẫu
        nhỏ.
 
  \item \textbf{Phân phối không cân bằng theo nhóm bệnh:}

        Nhóm bệnh nhân sa sút trí tuệ nặng (MMSE thấp) chiếm tỉ lệ nhỏ trong
        tập huấn luyện. Điều này giải thích hiện tượng quan sát được trên Q-Q
        Plot của tất cả các mô hình: sai số hệ thống tăng mạnh ở đuôi dưới phân
        phối—chính là các ca lâm sàng quan trọng nhất cần dự báo chính xác.
\end{itemize}

\newpage

\section{Bài học rút ra}

\begin{itemize}
  \item \textbf{Diagnostic plots đẹp trên tập Train không đồng nghĩa với mô
        hình tốt hơn.}

        Kernel Ridge có phần dư gần ngẫu nhiên và Scale-Location phẳng hơn
        Ridge trên tập huấn luyện, nhưng lại thua Ridge trên tập kiểm thử. Đây
        là bài học cốt lõi của Machine Learning: đánh giá mô hình phải luôn dựa
        trên tập dữ liệu \emph{chưa từng thấy}, không phải trên dữ liệu đã
        dùng để huấn luyện hay tinh chỉnh siêu tham số.
 
  \item \textbf{Độ phức tạp mô hình phải tương xứng với kích thước dữ liệu.}

        Với $n = 371$ và $p = 4$, Ridge (4 tham số $\hat{\beta}$) tổng quát hóa
        tốt hơn Kernel Ridge (296 dual coefficients). Tăng độ phức tạp mô hình
        không nhất thiết cải thiện hiệu năng—nguyên tắc \emph{Occam's Razor}
        áp dụng trực tiếp trong lựa chọn mô hình thực nghiệm.
 
  \item \textbf{Regularization và feature selection là hai triết lý khác nhau.}

        OLS chọn biến loại bỏ \texttt{EDUC} để đạt tính diễn giải tối giản
        (Parsimonious Model). Ridge giữ lại \texttt{EDUC} nhưng thu nhỏ hệ số
        thông qua phạt $L_2$—đây là sự đánh đổi có chủ đích giữa
        \emph{interpretability} và \emph{variance control}. Không có lựa chọn
        tuyệt đối đúng; quyết định phụ thuộc vào mục tiêu ứng dụng.
  
  \item \textbf{Kết quả $R^2 < 0$ là thông tin có giá trị khoa học,
        không phải thất bại.}

        Tất cả bốn mô hình đều ghi nhận $R^2_{\text{Test}} < 0$, khẳng định
        rằng tổ hợp tuyến tính của bốn đặc trưng hiện tại chưa đủ để phân tách
        tín hiệu suy giảm nhận thức mang tính cá thể hóa cao trên bộ dữ liệu
        này. Kết quả này định hướng rõ ràng cho nghiên cứu tiếp theo: cần thêm
        đặc trưng sinh học (\texttt{eTIV}, \texttt{CDR}), mô hình hóa
        interaction terms (\texttt{Age} $\times$ \texttt{nWBV}), hoặc tích hợp
        cấu trúc chuỗi thời gian của bộ dữ liệu longitudinal.
\end{itemize}
 